% HR Workflow Handbook — Aureus ERP (Truck It In (Private) Limited)
% Compile with pdfLaTeX on Overleaf (default compiler is fine).
%
% Structure:
%   Part I   — Workflow Reference   (how each HR process actually works today)
%   Part II  — Issues Found & Fixed (transparency log for this testing pass)
%   Part III — Questions for HR Review (kept separate from Part I, as requested)
%   Sign-off block

\documentclass[11pt,a4paper]{article}

\usepackage[margin=2.4cm]{geometry}
\usepackage[table]{xcolor}
\usepackage{enumitem}
\usepackage{longtable}
\usepackage{array}
\usepackage{booktabs}
\usepackage{titlesec}
\usepackage{fancyhdr}
\usepackage{parskip}
\usepackage{amssymb}
\usepackage{tcolorbox}
\tcbuselibrary{skins,breakable}
\usepackage[colorlinks=true,linkcolor=accentdark,urlcolor=accentdark]{hyperref}

% ---------- palette ----------
\definecolor{accent}{HTML}{2F5D50}
\definecolor{accentdark}{HTML}{1D3A32}
\definecolor{accentsoft}{HTML}{DBE8E2}
\definecolor{amberc}{HTML}{8A5A12}
\definecolor{ambersoft}{HTML}{F3E3C8}
\definecolor{inkc}{HTML}{1C2422}
\definecolor{inksoft}{HTML}{4D5A56}
\definecolor{linec}{HTML}{D7DDDA}
\definecolor{surface2}{HTML}{EEF1F0}

% ---------- section styling ----------
\titleformat{\section}
  {\large\bfseries\color{inkc}}{\thesection}{0.6em}{}
\titleformat{\subsection}
  {\bfseries\color{accentdark}}{}{0em}{}
\titlespacing*{\subsection}{0pt}{1.4em}{0.3em}

\pagestyle{fancy}
\fancyhf{}
\renewcommand{\headrulewidth}{0.4pt}
\fancyhead[L]{\small\color{inksoft} HR Workflow Handbook}
\fancyhead[R]{\small\color{inksoft} Aureus ERP \textperiodcentered\ Truck It In (Private) Limited}
\fancyfoot[C]{\small\color{inksoft} \thepage}

% ---------- helper boxes ----------
\newtcolorbox{statusbox}[1]{
  colback=accentsoft, colframe=accent, boxrule=0.4pt, arc=2pt,
  left=8pt, right=8pt, top=4pt, bottom=4pt, boxsep=1pt,
  fontupper=\small\bfseries\color{accentdark}, title={#1}, titlerule=0pt
}
\newtcolorbox{callout}{
  colback=accentsoft, colframe=linec, leftrule=3mm, colframe=accent,
  arc=1pt, boxrule=0.3pt, left=10pt, right=10pt, top=8pt, bottom=8pt
}
\newtcolorbox{note}{
  colback=surface2, colframe=linec, boxrule=0.4pt, arc=1pt,
  left=9pt, right=9pt, top=6pt, bottom=6pt
}
\newtcolorbox{questionsbox}[1]{
  colback=white, colframe=amberc, boxrule=0.6pt, arc=2pt, breakable,
  left=10pt, right=10pt, top=8pt, bottom=8pt,
  title={\color{amberc}\bfseries #1}, coltitle=amberc,
  fonttitle=\bfseries, colbacktitle=ambersoft, boxsep=2pt
}

\newcommand{\eyebrow}[1]{{\small\ttfamily\color{accentdark} #1}\par\vspace{2pt}}
\newcommand{\wftitle}[1]{{\Large\bfseries\color{inkc} #1}\par\vspace{6pt}}
\newcommand{\qitem}{\item[$\square$]}

\title{\vspace{-2em}{\bfseries\Huge HR Workflow Handbook}\\[6pt]
\large\color{inksoft} A step-by-step reference for every HR process in Aureus ERP,\\
plus the exact questions to walk through with HR before testing continues}
\author{}
\date{\small Company on file: Truck It In (Private) Limited \quad\textbar\quad
Environment: aureuserp @ 127.0.0.1:3307 \quad\textbar\quad Prepared: 7 Sep 2026}

\begin{document}
\maketitle
\thispagestyle{fancy}

\begin{callout}
\textbf{How to use this document.} Part~I explains how each HR process works today, step by step, in plain language. Part~II is a short transparency log of what we found broken and fixed this week. Part~III — kept deliberately separate — is the exact list of questions to walk through with HR, so her feedback comes back as concrete, actionable corrections rather than a general ``looks fine.''
\end{callout}

\tableofcontents

\newpage
\part*{Part I --- Workflow Reference}
\addcontentsline{toc}{part}{Part I --- Workflow Reference}

\section*{Before the workflows --- four ideas used everywhere}
\eyebrow{REFERENCE}
Four concepts repeat across almost every process below. Understanding them once makes the rest of this document much shorter to read.

\begin{enumerate}[label=\textbf{\arabic*.}, leftmargin=1.6em]
  \item \textbf{Company boundary.} Every record --- employee, leave, timesheet, candidate --- belongs to exactly one company. A person only ever sees records from their own company.
  \item \textbf{Login vs.\ HR record.} A \emph{User} is a login (email + password). An \emph{Employee} is the HR record (job title, salary, manager, leave balance). Most people have both, linked together --- but not every login needs an employee record, and not every employee needs a login.
  \item \textbf{Reporting hierarchy.} Who can see or approve for whom is decided by three things together: direct manager (and their manager, and so on), department, and team. Someone with the HR administrator override can see everyone regardless.
  \item \textbf{One shared approval engine.} Leave, timesheets, salary changes, and expense claims all route through the \emph{same} approval mechanism, just configured differently per request type. Learn it once in Time Off (\S 3) and every other approval works the same way.
\end{enumerate}

\section{Employee onboarding}
\eyebrow{\S01 \textperiodcentered\ People}
\begin{statusbox}{VERIFIED IN THIS ENVIRONMENT}\end{statusbox}

Creating a new employee record. This is separate from creating a login --- HR fills in the person's details, and the system quietly builds the supporting records around them.

\begin{enumerate}[leftmargin=1.8em]
  \item HR opens \textbf{Employees → New} and fills in name, job title, department, manager and company.
  \item \textbf{If the Company field is left blank} (common --- it sits on a secondary tab): the system fills it in automatically from the HR user's own company.
  \item \textbf{If it was set explicitly:} the chosen company is used as-is.
  \item On save, a matching \emph{Partner} contact record is created automatically and linked back to the employee.
  \item A first status-history entry is logged (e.g.\ ``Active,'' joining date).
\end{enumerate}

\begin{center}
\small
\begin{tabular}{@{}p{4cm}p{9.5cm}@{}}
\toprule
\textbf{Who does this} & HR administrator, or a manager with employee-management access. \\
\textbf{Automatic side effects} & Linked contact (Partner) record; status-history entry marking the employee active. \\
\textbf{Why the auto-fill matters} & Without it, a record with no company becomes invisible to every company-scoped screen --- including the confirmation page right after saving. \\
\bottomrule
\end{tabular}
\end{center}

\begin{note}
\textbf{Fixed this week:} the company auto-fill above (and the same fix for job positions, work locations, and performance records) was missing until 7 Sep --- see Part~II.
\end{note}

\section{Reporting lines \& who can see whom}
\eyebrow{\S02 \textperiodcentered\ Access}
\begin{statusbox}{VERIFIED IN THIS ENVIRONMENT}\end{statusbox}

Determines what any one person sees in Employees, Time Off, Timesheets, Performance and Attendance. Checked every time a list of employees or their records is shown.

\begin{enumerate}[leftmargin=1.8em]
  \item Person opens an HR screen.
  \item \textbf{If they hold the HR ``view all'' override:} they see every employee in their company --- skip to step 6.
  \item Otherwise, the system walks their manager chain downward (direct and indirect reports).
  \item It adds their department (and any sub-departments).
  \item It adds any team they manage. The union of steps 3--5 is what they can see.
  \item For any target record: if its company does not match the viewer's company, access is denied regardless of hierarchy; otherwise access is allowed.
\end{enumerate}

\begin{center}
\small
\begin{tabular}{@{}p{7.2cm}p{6.3cm}@{}}
\toprule
\textbf{Scenario} & \textbf{Result} \\
\midrule
Manager viewing a direct report & Allowed --- full record \\
Manager viewing an unrelated employee, same company & Not allowed unless same department/team \\
Anyone viewing an employee in a different company & Always blocked, even for HR admins without cross-company grant \\
HR admin with ``view all records'' permission & Sees everyone in their own company \\
\bottomrule
\end{tabular}
\end{center}

\begin{note}
Confirmed live: logged in as a manager with three visible reports, a fourth employee outside the hierarchy correctly did not appear.
\end{note}

\section{Time off \& leave}
\eyebrow{\S03 \textperiodcentered\ Leave}
\begin{statusbox}{VERIFIED IN THIS ENVIRONMENT}\end{statusbox}

Four stages: HR defines the kinds of leave that exist, HR allocates a balance to each employee, the employee requests time off against that balance, and the manager approves it.

\begin{enumerate}[leftmargin=1.8em]
  \item \textbf{HR (one-time setup):} create a Leave Type (e.g.\ Annual Leave), then create an Allocation (e.g.\ 14 days for one employee), then approve that allocation.
  \item \textbf{Employee:} goes to My Time Off → New, picks a Leave Type and a date range.
  \item \textbf{If the allocated balance is insufficient:} the request is blocked with a message and cannot be submitted.
  \item \textbf{If sufficient:} the request is created in state \texttt{confirm}.
  \item Employee clicks ``Submit for Approval''; the request routes to the requester's direct manager.
  \item \textbf{Manager approves:} state becomes \texttt{validate\_two} and the days are deducted from the balance.
  \item \textbf{Manager rejects:} state becomes \texttt{refuse} and the balance is left untouched.
\end{enumerate}

\begin{center}
\small
\begin{tabular}{@{}p{4.4cm}p{9.1cm}@{}}
\toprule
\textbf{Screen path (employee)} & Time Off → My Time → My Time Off → New \\
\textbf{Screen path (manager)} & Time Off → Management → Time Off → Waiting For Me \\
\textbf{Setup path (HR)} & Time Off → Configurations → Leave Types, and → Allocations \\
\textbf{One-time setup needed} & An approval routing rule for ``leave requests'' must exist per company \\
\bottomrule
\end{tabular}
\end{center}

\begin{note}
\textbf{Two things HR must set up once per company} before anyone can request leave: (1)~at least one active Leave Type, and (2)~an Approval Workflow for leave requests (Settings → Approval Workflows). Both were missing for this company until today --- see Part~II.
\end{note}

\section{Attendance}
\eyebrow{\S04 \textperiodcentered\ Time tracking}
\begin{statusbox}{VERIFIED BY AUTOMATED TEST --- NOT YET WALKED THROUGH ON SCREEN}\end{statusbox}

Records a single day's check-in/check-out against the employee's scheduled hours, and the system works out lateness and worked hours on its own.

\begin{enumerate}[leftmargin=1.8em]
  \item An attendance record is saved with the scheduled shift and the actual check-in/check-out times.
  \item The system computes worked hours, late minutes, and early-departure minutes automatically.
\end{enumerate}

\begin{center}
\small
\begin{tabular}{@{}lccc@{}}
\toprule
\textbf{Example} & \textbf{Scheduled} & \textbf{Actual} & \textbf{System computes} \\
\midrule
Sample day tested & 09:00--17:00 & 09:15--16:45 & 7.5 hrs worked; 15 min late; 15 min early departure \\
\bottomrule
\end{tabular}
\end{center}

\begin{note}
Confirmed by an automated test that reproduces exactly this scenario. Not yet clicked through on the live screens in this session --- worth a manual pass before relying on it for payroll-adjacent decisions.
\end{note}

\section{Performance reviews}
\eyebrow{\S05 \textperiodcentered\ Growth}
\begin{statusbox}{VERIFIED BY AUTOMATED TEST --- NOT YET WALKED THROUGH ON SCREEN}\end{statusbox}

A review cycle (e.g.\ ``2026 Annual Review'') is launched once and creates one review per employee automatically. Each review has two independent scores.

\begin{enumerate}[leftmargin=1.8em]
  \item HR launches a Performance Cycle with a start and end date.
  \item A Review is created for every active employee automatically.
  \item Employee submits a self-rating with comments.
  \item Manager submits a manager rating with comments.
  \item Review status becomes \texttt{completed}.
  \item \emph{Optional:} individual Goals can be attached to a review, each with its own weight and target value.
\end{enumerate}

\begin{center}
\small
\begin{tabular}{@{}lccc@{}}
\toprule
\textbf{Example tested} & \textbf{Self rating} & \textbf{Manager rating} & \textbf{Final status} \\
\midrule
2026 Annual Review, one employee & 4.1 / 5 & 4.4 / 5 & completed \\
\bottomrule
\end{tabular}
\end{center}

\section{Timesheets}
\eyebrow{\S06 \textperiodcentered\ Time tracking}
\begin{statusbox}{VERIFIED BY AUTOMATED TEST --- NOT YET WALKED THROUGH ON SCREEN}\end{statusbox}

A logged line of billable or internal work goes through the same submit-then-approve pattern as leave, and locks once approved so it can be relied on for billing.

\begin{enumerate}[leftmargin=1.8em]
  \item Employee logs hours; timesheet starts in state \texttt{draft}.
  \item Employee submits it for approval.
  \item \textbf{Manager approves:} state becomes \texttt{approved} and the entry locks --- auditable, recording who approved it and when.
  \item \textbf{Manager rejects:} state becomes \texttt{rejected} and it is sent back to the employee.
\end{enumerate}

\begin{note}
Every approval records who approved it and when --- the same underlying mechanism used for leave and salary changes (see the four ideas at the start of Part~I).
\end{note}

\section{Recruitment \& hiring}
\eyebrow{\S07 \textperiodcentered\ Hiring}
\begin{statusbox}{VERIFIED BY AUTOMATED TEST --- NOT YET WALKED THROUGH ON SCREEN}\end{statusbox}

From an open role to a hired employee, with no manual re-typing of the candidate's details at any step --- and no duplicate employee record if the same offer gets processed twice.

\begin{enumerate}[leftmargin=1.8em]
  \item A Job Position is published (e.g.\ company website, LinkedIn).
  \item A Candidate applies --- either entered manually, or pulled in automatically from a job board or API.
  \item An Applicant record tracks screening, interview and assessment scores.
  \item \textbf{If the offer status is ``accepted'':} the applicant is converted to an Employee.
  \item \textbf{If conversion has already run before for this applicant:} the existing employee is returned --- no duplicate is created.
  \item \textbf{If this is the first conversion:} a new Employee is created and the application is marked \texttt{hired}.
\end{enumerate}

\begin{center}
\small
\begin{tabular}{@{}p{4cm}p{9.5cm}@{}}
\toprule
\textbf{Duplicate protection} & Re-running the conversion, or re-importing the same external application ID, always resolves to the same one candidate/applicant/employee --- never a second copy. \\
\textbf{Company boundary} & A candidate imported for one company can never be attached to another company's job posting, even via the automated intake path. \\
\textbf{Reporting} & Headcount, funnel-by-stage, and average time-to-hire are available per company once hires start coming through. \\
\bottomrule
\end{tabular}
\end{center}

\section{Employee \& expense requests}
\eyebrow{\S08 \textperiodcentered\ Money}
\begin{statusbox}{VERIFIED BY AUTOMATED TEST --- NOT YET WALKED THROUGH ON SCREEN}\end{statusbox}

Where HR meets Accounting: an approved financial request (an expense claim, reimbursement, travel advance, etc.) automatically becomes a draft accounting entry --- never posted automatically, always left for Accounting to review.

\begin{enumerate}[leftmargin=1.8em]
  \item Employee files a request (e.g.\ ``Client travel reimbursement,'' PKR 5{,}000) and submits it for approval.
  \item \textbf{Finance/manager rejects:} state becomes \texttt{rejected}.
  \item \textbf{Finance/manager approves:} state becomes \texttt{approved}, and a balanced draft accounting journal is created automatically --- debit expense, credit payable.
  \item The journal sits in Accounting as a draft, awaiting Accounting's own posting.
\end{enumerate}

\begin{center}
\small
\begin{tabular}{@{}lcccc@{}}
\toprule
\textbf{Example tested} & \textbf{Amount} & \textbf{Debit} & \textbf{Credit} & \textbf{Journal state} \\
\midrule
Client travel reimbursement & PKR 5{,}000.00 & Expense account & Payable account & draft \\
\bottomrule
\end{tabular}
\end{center}

\begin{note}
Re-approving or re-processing the same request never creates a second journal entry --- only ever exactly one, kept in sync.
\end{note}

\section{Sensitive record changes}
\eyebrow{\S09 \textperiodcentered\ Compliance}
\begin{statusbox}{VERIFIED IN THIS ENVIRONMENT}\end{statusbox}

A short, fixed list of an employee's most sensitive fields can never be edited directly by anyone, even HR --- a change to any of them has to go through approval first, and the audit trail keeps a written reason.

\begin{enumerate}[leftmargin=1.8em]
  \item Someone proposes a change to a protected field; an approval request is created with both the old and new values recorded.
  \item \textbf{Approver rejects:} no change is made.
  \item \textbf{Approver approves (with a written reason):} the change is applied to the employee record, and the reason is permanently kept on the approval's audit trail.
\end{enumerate}

\begin{center}
\small
\begin{tabular}{@{}p{4cm}p{9.5cm}@{}}
\toprule
\textbf{Protected fields} & \texttt{identification\_id}, \texttt{passport\_id}, \texttt{ssnid}, \texttt{sinid}, \texttt{bank\_account\_id}, \texttt{salary\_grade}, \texttt{base\_salary}, \texttt{salary\_currency\_id} \\
\textbf{Anything not on that list} & Quietly ignored, even if bundled into the same request --- e.g.\ a name change submitted alongside a salary change is dropped; only the salary change goes through approval. \\
\bottomrule
\end{tabular}
\end{center}

\begin{note}
Known, deliberate scope limit: this protects the field from silent edits \emph{through the approval workflow}. It does not yet stop someone with direct database or admin-tooling access --- that is a separate, larger control and out of scope for this pass.
\end{note}

\newpage
\part*{Part II --- Issues Found \& Fixed This Week}
\addcontentsline{toc}{part}{Part II --- Issues Found \& Fixed This Week}

For transparency --- things that did not work correctly during testing, found and corrected in the same session. Included so HR is not the one who finds these in production.

\begin{longtable}{@{}p{2.6cm}p{5.6cm}p{5.6cm}c@{}}
\toprule
\textbf{Area} & \textbf{What was wrong} & \textbf{Fix} & \\
\midrule
\endhead
Employee onboarding &
Leaving the Company field blank on the create form (easy to miss --- it's on a second tab) made the new employee invisible everywhere, including the confirmation screen. &
Now defaults to the HR user's own company automatically. Same fix applied to job positions, work locations, and performance records. &
\cellcolor{accentsoft}\textcolor{accentdark}{\textbf{fixed}} \\
\addlinespace
Time off &
A Leave Type created through the normal admin screen had no way to be marked active --- so it silently never appeared in anyone's leave-request dropdown. &
Added the missing Active toggle to the Leave Type screen, defaulted new ones to active, and fixed the existing ``Annual Leave'' record. &
\cellcolor{accentsoft}\textcolor{accentdark}{\textbf{fixed}} \\
\addlinespace
Time off &
No approval routing existed for leave requests in this company at all, so ``Submit for Approval'' failed outright with a server error. &
Configured the approval routing (leave $\rightarrow$ requester's direct manager), and added ``leave request'' to the suggested list so it's easy to find next time. &
\cellcolor{accentsoft}\textcolor{accentdark}{\textbf{fixed}} \\
\bottomrule
\end{longtable}

\begin{note}
This is a running log for this testing pass --- earlier phases (invoicing, bank reconciliation, HR authorization) each closed with their own detailed reports, available on request.
\end{note}

\newpage
\part*{Part III --- Questions for HR Review}
\addcontentsline{toc}{part}{Part III --- Questions for HR Review}

\begin{callout}
These are kept separate from Part~I on purpose. Walk through each block below with HR \emph{after} she has read the matching section --- her answers tell us exactly what to change, and where.
\end{callout}

\begin{questionsbox}{\S01 --- Onboarding}
\begin{itemize}[leftmargin=1.6em]
  \qitem Who is actually allowed to create an employee record --- HR only, or can a manager add their own hire?
  \qitem Is there anything required before someone counts as ``active'' that the system doesn't ask for (signed contract, ID check, etc.)?
\end{itemize}
\end{questionsbox}

\begin{questionsbox}{\S02 --- Reporting lines \& visibility}
\begin{itemize}[leftmargin=1.6em]
  \qitem Should a manager see their whole department, or just people who report to them directly? (System currently does both --- manager chain \emph{and} department \emph{and} team.)
  \qitem Does HR need to see every company, or only the one they're assigned to?
\end{itemize}
\end{questionsbox}

\begin{questionsbox}{\S03 --- Time off}
\begin{itemize}[leftmargin=1.6em]
  \qitem Is manager approval the \emph{only} approval step, or does HR/a second person also need to sign off for leave? (System is single-step right now.)
  \qitem What leave types should actually exist (Annual, Sick, Casual, Unpaid\ldots) and how many days does each role get?
  \qitem Should HR be notified when someone requests leave, even if they're not the approver?
\end{itemize}
\end{questionsbox}

\begin{questionsbox}{\S04 --- Attendance}
\begin{itemize}[leftmargin=1.6em]
  \qitem Is 15 minutes late actually ``late,'' or is there a grace period (e.g.\ 10 minutes) before it counts?
\end{itemize}
\end{questionsbox}

\begin{questionsbox}{\S05 --- Performance}
\begin{itemize}[leftmargin=1.6em]
  \qitem How often do review cycles actually run --- annual, twice a year?
  \qitem Is self-rating + manager-rating the whole process, or is there a final HR calibration step missing?
\end{itemize}
\end{questionsbox}

\begin{questionsbox}{\S06 --- Timesheets}
\begin{itemize}[leftmargin=1.6em]
  \qitem Is manager approval enough, or does Finance also sign off on billable hours before invoicing?
\end{itemize}
\end{questionsbox}

\begin{questionsbox}{\S07 --- Recruitment}
\begin{itemize}[leftmargin=1.6em]
  \qitem Does an accepted offer need budget/Finance sign-off, or is Recruitment/HR's decision final?
\end{itemize}
\end{questionsbox}

\begin{questionsbox}{\S08 --- Expense requests}
\begin{itemize}[leftmargin=1.6em]
  \qitem What expense categories should exist, and should large amounts route to someone above the direct manager?
\end{itemize}
\end{questionsbox}

\begin{questionsbox}{\S09 --- Sensitive changes}
\begin{itemize}[leftmargin=1.6em]
  \qitem Is that protected-field list complete? Anything missing, like job title or termination date?
  \qitem Should salary changes specifically require Finance approval, not just any manager?
\end{itemize}
\end{questionsbox}

\vspace{1.5em}
\noindent\textbf{Reporting back:} you don't need to translate her answers into system language --- just note, per section, what she said is different from what's described (e.g.\ \emph{``\S03 --- leave needs two approvals: manager then HR, not just manager''}, or \emph{``\S08 --- anything over PKR 20{,}000 needs the Finance Manager, not the direct manager''}). Each answer becomes a concrete configuration or code change.

\newpage
\section*{Review \& sign-off}
\addcontentsline{toc}{section}{Review \& sign-off}

\noindent
\begin{tabular}{@{}p{3.4cm}p{10cm}@{}}
\textbf{Reviewed by} & \rule{9cm}{0.4pt} \\[1.6em]
\textbf{Role} & \rule{9cm}{0.4pt} \\[1.6em]
\textbf{Date} & \rule{9cm}{0.4pt} \\[1.6em]
\textbf{Sections needing changes} & \rule{9cm}{0.4pt} \\[1.6em]
 & \rule{9cm}{0.4pt} \\[1.6em]
 & \rule{9cm}{0.4pt} \\
\end{tabular}

\vspace{2em}
\begin{center}
\small\color{inksoft}
Aureus ERP \textperiodcentered\ HR Workflow Handbook \textperiodcentered\ generated from the live test environment, 7 Sep 2026
\end{center}

\end{document}
